\section{Bug Bounty y Safe Harbor: la recompensa no sustituye los límites de la autorización}

Un VDP y un bug bounty suelen aparecer juntos, pero resuelven problemas distintos. Un \term{bug bounty program} (programa de recompensas por vulnerabilidades) es un programa opcional que, bajo reglas establecidas, ofrece dinero, puntos o reconocimiento a los reportes válidos que cumplen las condiciones; un VDP debe existir aunque no haya recompensa. El bounty ayuda a competir por la atención de los investigadores, a fijar severity rewards y a operar el triage, pero un pago posterior no puede reescribir una conducta dañina que ya ocurrió.

\subsection{El límite entre VDP y bounty}

Una organización puede tener un VDP sin bounty, y también puede administrar el bounty mediante una plataforma de terceros. Ambos comparten el alcance de activos, el canal de reporte y el mecanismo de coordinación, pero el bounty requiere además reward table, duplicate rule, eligibility, tax/payment y appeal. Los responsables de seguridad deben evitar que el equipo pase por alto data access, extortion signal u otros incident indicators porque «es un reporte de bounty».

\lecturefigure{04-vdp-vs-bounty.png}{El VDP ofrece un canal seguro de reporte; el bug bounty agrega sobre él un mecanismo opcional de recompensa.}{CISA/DOJ VDP guidance; subtítulos locales 12:30--15:20; redibujado conceptual.}

\begin{knowledgebox}{Lectura de la figura: primero la autorización, después la recompensa}
El núcleo del VDP es safe channel, scope, response promise y coordinated disclosure; la capa adicional del bounty es optional payment, severity table, eligibility y duplicate. Ninguno de los dos cubre la extortion, el sabotaje, el acceso a datos más allá de lo necesario ni la negativa a detener las pruebas. Al evaluar un reporte, primero se determina si la conducta se ajustó a los límites publicados de antemano y a la good faith, y después se valoran la vulnerabilidad y la recompensa; el proceso de pago no debe decidir la clasificación legal ni la del incidente.
\end{knowledgebox}

\begin{warningbox}{El payment no es una post-hoc authorization}
Que la organización pague una recompensa, firme un NDA o mantenga la comunicación puede servir para controlar el riesgo y facilitar la corrección; estas acciones no prueban por sí mismas que la conducta previa estuviera autorizada, ni eliminan la necesidad de conservar evidencia, notificar a los afectados o consultar a counsel. Cualquier atajo del tipo «pagar para convertir el incident en bounty» crea un punto ciego de gobernanza.
\end{warningbox}

\subsection{Safe harbor, scope y good faith}

En el contexto de un VDP, el \term{safe harbor} (puerto seguro) es el compromiso de la organización con la investigación de buena fe que cumple la política; por ejemplo, no ejercer ciertas acciones civiles, respaldar al investigador para que demuestre su buena fe o, dentro del ámbito aplicable, no recomendar acciones penales. No es una inmunidad legal general ni puede obligar a todos los terceros. El \term{scope} (alcance) especifica los dominios, aplicaciones, API, versiones y activos excluidos que pueden probarse; la \term{good faith} (buena fe) suele exigir solo el acceso mínimo necesario para la verificación, evitar afectaciones a la privacidad y al servicio, y detenerse y reportar a tiempo.

\lecturefigure{05-safe-harbor.png}{El safe harbor solo es eficaz si se combina con un scope claro, requisitos de buena fe, conductas prohibidas y una ruta de escalamiento.}{CISA VDP template; DOJ VDP framework; redibujado conceptual.}

\begin{knowledgebox}{Lectura de la figura: la política debe proteger a la vez a los investigadores y a los usuarios}
In-scope enumera los activos y métodos que se permite probar; good faith exige minimizar el acceso, reportar y proteger los datos sensibles; prohibited fija límites como disruption, persistence, data use o extortion; escalation ofrece un safety contact, un legal contact y una vía para las disputas. Si solo hay frases de aliento sin un scope concreto, el investigador asume riesgos impredecibles; si solo hay una lista de prohibiciones sin compromiso de respuesta, no se genera confianza.
\end{knowledgebox}

\begin{importantbox}{Lista de verificación para diseñar la política}
El inventario de activos debe poder mantenerse; los servicios de terceros deben indicar los límites de propiedad; las restricciones de velocidad de prueba y de automatización deben explicar su motivo; el manejo de datos sensibles debe establecer cuándo detenerse, la retención mínima y la eliminación; las disputas deben contar con un escalation humano. Las actualizaciones de la política también deben conservar la versión y la fecha de entrada en vigor, para que la organización no evalúe conductas anteriores con reglas nuevas.
\end{importantbox}

\subsection{Resumen de la sección}

El VDP se ocupa del canal, el bounty de los incentivos y el safe harbor del compromiso de la organización con quien respeta los límites. Ninguno de los tres sustituye la incident classification, y ninguno permite tratar el pago, el NDA o la presunción de buena fe como hechos técnicos.
